% 自动生成，勿手改（gen_notes.py）
\chapter{译校批注}

本附录汇集正文中以源码注释形式给出的译校批注，按章节顺序排列。每条注明所在原书页码，便于对照原书。批注只写三类内容：易错预警（符号、指标、约定）、作者省略论证的点拨、背景直觉。原书疑似笔误一律照原样排，仅在批注中指出。\par\medskip

\section*{第 1 章\quad 什么是镜像对称}

\noindent （原书 p.11）\ \textbf{§ 镜像变换与有理曲线的条数}\ 复结构模参数在原文里记号并不统一：p.11 的叙述用 $\psi$，式（26）写作 $\phi$，下一页的式（27）又回到 $\psi$，三者是同一个参数。译文照录原记号，便于与原书对照。\par\smallskip

\section*{第 2 章\quad Gromov--Witten 不变量入门}

\noindent （原书 p.20）\ \textbf{§ 相交数}\ 原文此处写“（5）是（4）的特殊情形”，但（4）是相交数 $\#(P\cap Q)$，（3）才是 $P,Q$ 的一般方程组；按文意疑应为（3）。此处照原文录为（4）。\par\smallskip

\noindent （原书 p.22）\ \textbf{§ 同调}\ 本页黑体断言作「2 次同调群」，而脚注 4 中同一对象却印作「1 次」。$u_2$ 表示直线，对应的同调是实 2 维的，脚注的「1 次」疑为笔误；两处均照原文录出。\par\smallskip

\noindent （原书 p.28）\ \textbf{§ 计数几何}\ 原文此处印作 $d^1$。按除子公理，插入一个直线条件相当于乘上该曲线的次数 $d$；一般地，插入 $l$ 个这样的条件就给出因子 $d^l$。此处是 $l=1$ 的情形，原文照录。\par\smallskip

\noindent （原书 p.30）\ \textbf{§ Calabi--Yau 流形}\ 原文说（30）给出 $md+1$ 个条件式。每个系数取复数值，所以实的条件数是 $2(md+1)$——下文的减法正由此而来。\par\smallskip

\noindent （原书 p.31）\ \textbf{§ Calabi--Yau 流形}\ 原文此处印作「有現変数」，从上下文看疑为「有限変数」（有限个变量）之误，译文按文意处理。\par\smallskip

\section*{第 3 章\quad 正多面体与镜像对称}

\noindent （原书 p.33）\ \textbf{§ 引言}\ 问题 1 的正整数解分成无穷族 $(m,l-m+1,1)$ 与有限个例外 $(3,3,2)$、$(4,3,2)$、$(5,3,2)$；后三个正是正四面体、正八面体、正二十面体对应的三元组。它与 ADE 型单奇点、$SO(3)$ 有限子群、Dynkin 图 $A$--$D$--$E$ 的分类是同一件事。\par\smallskip

\noindent （原书 p.35）\ \textbf{§ 正多面体的分类}\ 所谓“塌缩的多面体”，就是二面体（dihedron）：把正 $p$ 边形的两面粘合，或把 $q$ 个“正 2 边形”沿边粘合。它的旋转对称群是二面体群，正好补上了后面 $SO(3)$ 有限子群分类中漏掉的那一族。\par\smallskip

\noindent （原书 p.36）\ \textbf{§ 正多面体群}\ 原文如此。$SO(3)$ 有限子群的完全分类是：循环群 $C_l$、二面体群 $D_l$、正四面体群 $T$、正八面体群 $O$、正二十面体群 $I$。原文把二面体群归入了上一段所说的“塌缩的多面体”的正多面体群，对照（6）式中指数取 $(2,2,q)$ 的情形即可理解。\par\smallskip

\noindent （原书 p.37）\ \textbf{§ 正多面体群}\ 图 1 中两类轴：过顶点与中心（对面中心）的轴，绕其旋转 $120^\circ$（对应 $b$）；过边中点与中心的轴，绕其旋转 $180^\circ$（对应 $c$）。正四面体的旋转对称群 $T$ 有 12 个元素，正是由这两种旋转生成的。\par\smallskip

\noindent （原书 p.37）\ \textbf{§ 正多面体群}\ （5）式即正四面体群 $T\cong A_4$ 的一组生成元与关系式。三个生成元对应三种旋转：$a$、$b$ 各为 $120^\circ$ 顶点旋转，$c$ 为 $180^\circ$ 边中点旋转；关系 $a^3=b^3=c^2=abc=1$ 的指数 $(3,3,2)$ 就是正四面体的 $(p,q,r)$。\par\smallskip

\noindent （原书 p.39）\ \textbf{§ 轨道空间的代数描述}\ 图 2 是“$1/p$ 点”的直观示意：左图中张角 $2\pi/p$ 的扇形把两条边粘合起来，就得到该尖点；右图中圆锥的顶点表示尖点，球面上被固定点的轨道被压成一个点。\par\smallskip

\noindent （原书 p.40）\ \textbf{§ 孤立超曲面奇点与轨道空间}\ 原文用 $\mathcal{C}_f$ 记带尖点信息的轨道空间，即商叠 $[X_f\setminus\{0\}/\C^*]$；用 $C_f$ 记它的粗模空间，后者是光滑紧黎曼面。\par\smallskip

\noindent （原书 p.42）\ \textbf{§ 与镜像对称的关联}\ 左栏与右栏正是镜像对称两侧的对应物：中间同调群对 Grothendieck 群，相交形式对 Euler 形式，Milnor 单值性对 Serre 函子。另，原书左栏印 $H_2(X_1,\Z)$，按上下文（$f=0$ 的纤维）疑为 $X_f$，此处照录底本。至于 p.41 的「D$\cdot$E 型奇点」，那是原书原样写法，指 ADE 型中的 D、E 两族。\par\smallskip

\section*{第 4 章\quad Calabi--Yau 流形与代数几何}

\noindent （原书 p.49）\ \textbf{§ 代数簇的分类与 Calabi--Yau 流形}\ 丰富性定理断言，3 维极小模型 $X$ 的典范除子 $K_X$ 是半丰沛的，即某个倍数 $|(K_X)^{\otimes n}|$ 无基点，从而 $\phi_n$ 是正则映射而不是仅仅有定义的映射。\par\smallskip

\noindent （原书 p.52）\ \textbf{§ K3 曲面}\ 本段开头“把 $S$ 的复结构稍作形变”一句，原书如此；严格说指在模空间中形变到 Kummer 曲面所在的开稠密集，未必是任意小的形变。\par\smallskip

\noindent （原书 p.53）\ \textbf{§ 3 维 Calabi--Yau 流形}\ 原文写作 $D:=\{(\zeta^i,\dots,\zeta^i)\in G\}$，但 $G$ 要到下文 $G:=\widetilde{G}/D$ 才定义，此处应为 $D\subset H$ 的笔误，译文照录原文。\par\smallskip

\noindent （原书 p.55）\ \textbf{§ 3 维 Calabi--Yau 流形}\ 这里把 $h^{1,1}(X)$ 定义为 $\dim H^1(X,\Omega^1_X)$，与 Hodge 理论中 $h^{1,1}=\dim H^1(X,\Omega^1)$ 的记法一致；对 Calabi--Yau 流形，$h^{1,1}$ 也是 Kähler 模的维数。\par\smallskip

\noindent （原书 p.55）\ \textbf{§ 3 维 Calabi--Yau 流形}\ 形变不改变拓扑，所以 Euler 数是形变不变量。小平定理说所有 $K3$ 曲面都能通过形变连起来。3 维 Calabi--Yau 流形的 Euler 数可以是 $-200$、$0$ 等等，彼此不同，因此无法只靠形变全部相连。这正是需要引入锥形流形跃迁的原因。\par\smallskip

\noindent （原书 p.57）\ \textbf{§ 锥形流形与 Reid 的幻想（Reid's fantasy）}\ 原书此页两幅图没有图题，图 1、图 2 的编号与题注为译者所加。图 1 中 $\tilde{V}$ 内的网格示意例外除子 $\PP^1\times\PP^1$；图 2 中的「変形」即“形变”，箭头表示锥形流形跃迁的第二步。\par\smallskip

\noindent （原书 p.58）\ \textbf{§ 锥形流形与 Reid 的幻想（Reid's fantasy）}\ 这就是所谓的 Reid 的幻想：把 3 维 Calabi--Yau 流形的“形变”与“带奇点的 Calabi--Yau 跃迁”两种操作反复施行，就能把任意两个 3 维单连通 Calabi--Yau 流形连起来。若为真，则 3 维单连通 Calabi--Yau 流形的双有理等价类可以连成一张“网”。\par\smallskip

\section*{第 5 章\quad 奇点与拓扑场论}

\noindent （原书 p.60）\ \textbf{§ 引言}\ 此句是说：非交换几何中的空间本来“光滑”，一旦用交换几何（复代数簇）来观察，其非交换性便表现为奇点。这里“非交换流形”是非交换几何中的空间，并非字面上的“非交换的流形”。\par\smallskip

\noindent （原书 p.63）\ \textbf{§ 量子上同调环的复习}\ 原书此处印作 $\delta_{i,n-i}$，下标只用到 $i$；但 $\eta_{ij}=C_{0ij}$ 是 $H^{2i}$ 与 $H^{2j}$ 的相交配对，非零当且仅当 $i+j=n$，故应为 $\delta_{i+j,n}$（等价于 $\delta_{j,n-i}$）。此处照录原书，读者按相交配对理解即可。\par\smallskip

\noindent （原书 p.63）\ \textbf{§ 量子上同调环的复习}\ 原文「（コ）ホモロジー」把「同调」与「上同调」并排写出，带有玩笑意味；正文译作「（上）同调」。\par\smallskip

\noindent （原书 p.64）\ \textbf{§ 量子上同调环的复习}\ 原书式（10）与 p.65 正文都印作 $t^i$；按坐标分量的写法与下文 $\deg t_i=1-i$，此处应为 $t_i$。为存原书面貌仍照录 $t^i$，读者按 $t_i$ 理解即可。\par\smallskip

\noindent （原书 p.65）\ \textbf{§ 量子上同调环的复习}\ 本段开头把 Chern 类直观地说成“衡量流形偏离仿射空间有多远的指标”，原书如此；严格说，Chern 类是切丛的示性类，第一 Chern 类为零正是 Calabi--Yau 条件。\par\smallskip

\noindent （原书 p.69）\ \textbf{§ 奇点、原始形式与平坦结构}\ 原书式（19）第一项印作 $t_0^2t_1$；按 $\deg t_0=1$、$\deg t_1=3/4$、$\deg t_2=1/2$，该项次数为 $11/4$，与其余各项的 $5/2$ 不符（$2\cdot1+1/2=5/2$，与 $EF_0=(5/2)F_0$ 一致的是 $t_0^2t_2$）。此处照录原书。\par\smallskip

\noindent （原书 p.70）\ \textbf{§ 结语}\ 「単純特異点」在此译作「单奇点」，指 A--D--E 型超曲面奇点（即二维 Du Val 奇点），与「单纯（simplicial）」无关；下文 Klein 正多面体的有限群与 Brieskorn 的根系描述正是指这一族奇点。\par\smallskip

\section*{第 6 章\quad 镜像对称与共形场论}

\noindent （原书 p.73）\ \textbf{§ 什么是弦理论}\ 原书此处写作 $R(G)_{i\bar j}=\partial_i\partial_{\bar j}\det(G_{m\bar n})$，照录原文；通常 Ricci 形式写作 $R_{i\bar j}=-\partial_i\partial_{\bar j}\log\det(G_{m\bar n})$。\par\smallskip

\noindent （原书 p.75）\ \textbf{§ 形变的参数与镜像对称}\ 「ミラー写像」译作「镜像映射」，与「ミラー変換」（镜像变换）分工：后者指两边的变换关系，前者指具体的那组映射。\par\smallskip

\noindent （原书 p.75）\ \textbf{§ 实 2 维环面的例子}\ 原书两处均写作“第 1 节”，但下文所指的“像是定义的东西”实际出现在第 2 节《什么是弦理论》中；这是底本的指代疏漏，译文照录不改。\par\smallskip

\noindent （原书 p.78）\ \textbf{§ 镜像对称与 T 对偶性}\ 此处原书写作 $\R/2\pi\Lambda_{\R}$，p.76 则写作 $\R^2/2\pi\Lambda_{\Z}$；两处均照录，疑为笔误。\par\smallskip

\noindent （原书 p.80）\ \textbf{§ 镜像对称与 T 对偶性}\ T 对偶是弦论的对称性：半径 $R$ 的圆周与半径 $1/R$ 的圆周给出同一个弦理论（此处取 $\alpha'=1$），故 $R_1\mapsto 1/R_1$ 不改变物理。\par\smallskip

\section*{第 7 章\quad 环簇几何与 Calabi--Yau 流形}

\noindent （原书 p.87）\ \textbf{§ 什么是环簇}\ 下文（7）式原文写作 $(z_1,z_2,z_3)\in(\C^\times)^3$，但紧接着的例子 $Z(\sigma_{\{1,2\}})=\{(z_1,z_2,0)\}$ 含有 $z_3=0$，可见原文的 $(\C^\times)^3$ 有误，当作 $\C^3$。此处照录原文，不改动。\par\smallskip

\noindent （原书 p.88）\ \textbf{§ 作为辛流形}\ 其一，$\CP^2/(\Z/3\Z)$ 有不动点，这里取的是轨形（orbifold）意义下的商；相应地 $X(\Sigma^\circ)$ 作为簇带有商奇点，与第 4 节的反射多面体构造正好接上。其二，上文把辛结构写成“非退化的 2 次微分形式”，非退化到下一句才解释；原书此处未提闭性 $\mathrm{d}\omega=0$，通常的辛形式定义还要求闭，译文照录。\par\smallskip

\noindent （原书 p.91）\ \textbf{§ 镜像对称}\ （15）的两式把 Hodge 数 $h^{1,1}$ 与 $h^{2,1}$ 互换：辛几何一侧的 Kähler 形变参数与复几何一侧的复结构形变参数维数相等。这正是「镜像」一词的拓扑含义，也是检查一个候选镜像对是否合格的第一道门槛。\par\smallskip

\noindent （原书 p.92）\ \textbf{§ 镜像对称}\ 此处第一个 $M^\circ$ 疑为 $M$，译文照录。\par\smallskip

\noindent （原书 p.93）\ \textbf{§ SYZ 猜想}\ “T 对偶”指半径 $R$ 的圆周与半径 $1/R$ 的圆周给出同一个物理的理论，是弦论中最早被认识的等价性；SYZ 猜想把镜像对称解释为沿纤维环面的 T 对偶。第 6 章会从共形场论的角度再讲这件事。\par\smallskip

\section*{第 8 章\quad Floer 同调与镜像对称}

\noindent （原书 p.95）\ \textbf{§ 引言}\ 原书此句印作「1 次元の図形」，但由 1 点构成的环路作为图形是 0 维的（正对应 $H_0$），原文此处自相矛盾，疑「1 次元」为「0 次元」之误；译文照原文录。\par\smallskip

\noindent （原书 p.99）\ \textbf{§ 佐藤 Grassmann 流形}\ 原书此处第二个 $\operatorname{Coker}$ 前似脱 $\dim$，按指标定理应为 $\dim\operatorname{Ker}(\pi|_L)-\dim\operatorname{Coker}(\pi|_L)=0$；且原书 $\operatorname{CoKer}$ 与 $\operatorname{Coker}$ 两处写法混印，此处均照原文录。\par\smallskip

\noindent （原书 p.100）\ \textbf{§ Hamilton 形式与 Lagrange 形式}\ 原书被积式中的 $\frac{q(t)}{dt}$ 按通行记号应作 $\frac{dq(t)}{dt}$，下文 $\mathfrak{H}$ 的式子同此；且 $\mathfrak{H}$ 的被积式印作 $p\frac{dq}{dt}+V$，与通行的 Hamilton 作用量 $p\frac{dq}{dt}-H$（此处 $H=\frac{p^{2}}{2}+V$）不同。均照原文录。\par\smallskip

\noindent （原书 p.100）\ \textbf{§ Hamilton 形式与 Lagrange 形式}\ 原书此处印作 $P(t)$，与前面定义的动量 $p(t)$ 大小写不一致；照原文录。\par\smallskip

\noindent （原书 p.102）\ \textbf{§ Floer 同调}\ 原书此处作“未知的”，按上下文疑为“闭的”之误；照原文录。\par\smallskip

\noindent （原书 p.105）\ \textbf{§ Floer 同调与变分问题}\ 原书此处引作（4），按上下文应为式（5）（p.103 中 $\mathfrak{H}$ 的定义式）；式（4）是 $\bigoplus_{k}H_{d-k}(\operatorname{Crit}_k(\mathfrak{L}))$。疑原书式号误，照原文录。\par\smallskip

\noindent （原书 p.105）\ \textbf{§ Floer 同调与变分问题}\ 原书（9）的被积式同样印作 $\frac{q(t)}{dt}$；按通行记号应作 $\frac{dq(t)}{dt}$，与第 4 节 $\mathfrak{L}$ 的写法同此，照原文录。\par\smallskip

\noindent （原书 p.105）\ \textbf{§ Floer 同调与变分问题}\ 原书此处印作 $t=s^{2\pi\sqrt{-1}s}$，原书如此；按上下文（$t\in S^{1}$ 的参数化）疑为 $t=e^{2\pi\sqrt{-1}s}$ 之误，正文照录。\par\smallskip

\noindent （原书 p.107）\ \textbf{§ Floer 同调与 Gromov--Witten 不变量}\ 原书此处印作「4 節の(5)」；式（5）在第 6 节且为 $\mathfrak{H}$，而 $\mathfrak{L}$ 在第 4、5 节中没有编号式。疑原书节号与式号均有误，照原文录。\par\smallskip

\noindent （原书 p.108）\ \textbf{§ Floer 同调与 Gromov--Witten 不变量}\ 原书此句作「$HF(M)$ 上のコホモロジー」，字面照录；$HF(M)$ 本身是 Floer 同调群，作者此说或指通过 (11) 在其上读出的量子上同调结构。\par\smallskip

\noindent （原书 p.108）\ \textbf{§ Floer 同调与 Gromov--Witten 不变量}\ 原书章末附有作者姓名的假名读音标记；译文按全书体例省略，作者名已在章首署名中给出。\par\smallskip

\section*{第 9 章\quad 镜像对称与 D-模}

\noindent （原书 p.113）\ \textbf{§ D-模}\ 此处原文作「記号的分解」，译文按中文文献通行的「Frobenius 分解」译出。\par\smallskip

\noindent （原书 p.114）\ \textbf{§ 以二维环面为例}\ $\CP^2$ 内“3 次超曲面”是余维 1 的对象，即一条 3 次曲线；原书沿用“超曲面”的称呼，译文照录。\par\smallskip

\noindent （原书 p.118）\ \textbf{§ 提取 Hodge 结构}\ VHS 是 variation of Hodge structure（变动的 Hodge 结构）的缩写。Hodge 数、Hodge 分解这些概念第 1 章与第 4 章已经用过；把它们随参数变动的整体情形一并处理，就叫 VHS。\par\smallskip

\noindent （原书 p.119）\ \textbf{§ 结语}\ 此处的 $J$ 是椭圆曲线模函数 $j$ 的 $1/1728$ 倍；$j=q^{-1}+744+196884q+\cdots$ 是模不变量，书中未加说明。\par\smallskip

\section*{第 10 章\quad 镜像对称与同调代数}

\noindent （原书 p.127）\ \textbf{§ 导出范畴}\ 原书把拟同构（quasi-isomorphism）称为「弱同构」，把导出范畴的对象称为「弱等价类」；以弱等价作局部化，正是导出范畴的标准造法。\par\smallskip

\noindent （原书 p.127）\ \textbf{§ 导出范畴}\ 原书此处作「两个向量丛 $E_1,E_2$」，而映射写作 $\partial:E_1\to E_0$，下文同调群亦用 $E_0$；疑「$E_2$」为「$E_0$」之误，译文照录原文。\par\smallskip

\noindent （原书 p.130）\ \textbf{§ A 模型与 B 模型}\ 原书 $\mathcal{T}$ 的指数印作 $\beta\cap\omega$；按上下文应理解为曲线类 $\beta$ 与 Kähler 形式 $\omega$ 的配对 $\int_\beta\omega$，即面积因子。\par\smallskip

\noindent （原书 p.133）\ \textbf{§ Floer 同调与同调镜像对称}\ 原书此页两幅图以「図 1」「図 2」标识，本译本的图号与题注为译者所加。图 1 的阴影表示以 $L_0,L_1$ 的弧为边界的全纯圆盘；图 2 的阴影表示由 $L_0,L_1,L_2$ 围成的全纯三角形。\par\smallskip

\noindent （原书 p.135）\ \textbf{§ 同调代数与场论}\ BRST 反常（anomaly）指经典对称性在量子化后不能保持的现象；此处指（1）式 $d_{k-1}\circ d_k=0$ 在一般 Floer 同调中只能在高阶（$A_\infty$）意义下成立。\par\smallskip

\noindent （原书 p.136）\ \textbf{§ 同调代数与场论}\ 重整化在这里指把中间状态（高能自由度）的效应吸收进低能有效理论的做法；正是它把原本只到 3 次的积结构提升为含 4 次以上高阶项的 $A_\infty$ 结构。\par\smallskip

\noindent （原书 p.136）\ \textbf{§ 同调代数与场论}\ 原书本章末尾附有作者姓名的假名读音标记，译文按全书体例省略；作者名已在章首署名中给出。\par\smallskip

\section*{第 11 章\quad 奇点与镜像对称}

\noindent （原书 p.138）\ \textbf{§ Hodge 数与拓扑镜像对称}\ 原书此处 $f_k$ 的变量印作 $x_1,\dots,x_n$，与前面的 $f_1(x_0,\dots,x_n)$ 不一致；按 $X\subset\CP^n$ 的齐次坐标应为 $x_0,\dots,x_n$。原书如此，照录。\par\smallskip

\noindent （原书 p.139）\ \textbf{§ Jacobi 环的 Poincaré 多项式}\ 原书此式印作 $\frac{1}{l+1}\cdot\frac{\partial f}{\partial x}=x^{l+1}=f(x)$。按定义 $Ef=\frac{a}{h}x\frac{\partial f}{\partial x}$，应为 $\frac{1}{l+1}\cdot x\cdot\frac{\partial f}{\partial x}=x^{l+1}=f(x)$，原书漏印了一个 $x$。照录原书。\par\smallskip

\noindent （原书 p.141）\ \textbf{§ 不可思议的相似性}\ 式（4）中印作 $(p-n/2)^2$，同页其余处一律用 $dx=\dim_{\C}X$；二者指同一量，属原书记号不统一。照录原文。\par\smallskip

\noindent （原书 p.142）\ \textbf{§ 轨形 Poincaré 多项式}\ 原书此处印作 $\omega_i:=\alpha_i/h$ 以及 $\alpha_i\in\Z$；但 $\alpha_i$ 同时又是群元素的分量（Vafa 公式中 $\alpha,\beta\in G$ 求和），按第 5 章式（2）的权 $a_i$，此处疑应为 $\omega_i:=a_i/h$。照录原书。\par\smallskip

\noindent （原书 p.144）\ \textbf{§ Milnor 纤维的几何学}\ 原书此处作「原始 $l+1$ 次単位根」，即本原（primitive）单位根；「原始」在本译本中固定对应 primitive form（原始形式），故正文改用「本原」以免混淆。\par\smallskip

\section*{第 12 章\quad 超越镜像对称}

\noindent （原书 p.151）\ \textbf{§ 局部镜像对称}\ 原文「動かない」直译为“不动”，指该 $\PP^{1}$ 的法丛没有全纯截面、曲线不能变形，义近几何中的“刚性”（rigid）。\par\smallskip

\noindent （原书 p.152）\ \textbf{§ 局部镜像对称}\ 原文「幾何工学」是 geometric engineering 的日译，中文文献通行作“几何工程学”，此处照此译出。\par\smallskip

\noindent （原书 p.153）\ \textbf{§ 大 N 相对性}\ 原书此处作「ラージ N 相対性」。按内容（Gopakumar--Vafa 1998 年的结果）通常称为 large N duality，即大 N 对偶；此处照原题译作“大 N 相对性”。\par\smallskip

\noindent （原书 p.158）\ \textbf{§ 几何工程——geometric engineering}\ 原书此处作「$G$-半自己双対接続」，通常称为 $G$-反自对偶联络（$G$-instanton），照录原词并说明。\par\smallskip

\noindent （原书 p.160）\ \textbf{§ Donaldson--Thomas 不变量}\ 这里 $\cO_X/I$ 是曲线 $f(C)$ 的结构层，故 $\chi(\cO_X/I)=\chi(\cO_C)=1-g$；亏格 $g$ 的信息就这样被 Euler 示性数吸收。\par\smallskip

\noindent （原书 p.161）\ \textbf{§ Donaldson--Thomas 不变量的跨墙公式}\ 这里的“导出范畴”即 derived category 的中译；原书第 10 章曾用「層の圏」与「導来圏」对举，读者可回看该章。\par\smallskip

\noindent （原书 p.162）\ \textbf{§ Donaldson--Thomas 不变量的跨墙公式}\ $\mathrm{Stab}(D(X))$ 是 Bridgeland 稳定性条件的参数空间（复流形），其上的“墙”是若干实余维 1 的壁；跨墙公式即比较墙两侧的不变量。\par\smallskip

\noindent （原书 p.163）\ \textbf{§ Donaldson--Thomas 不变量的跨墙公式}\ $b_+$ 是 4 维流形上相交形式正定部分的维数；$b_+=1$ 时 Donaldson 理论的模空间紧化可控，Kotschick--Morgan 猜想才有着落。\par\smallskip
